\chapter{Measuring Success: Edit Distance, WER, CER and LID Accuracy}
\chaptermark{Metrics}
\label{ch:metrics}

Before building anything we must agree on how to score it. A wrong metric would make every later comparison meaningless, and two
metric bugs in this project (Chapter~\ref{ch:lessons}) would have changed conclusions.

\section{Edit distance}

\begin{definition}{Levenshtein distance}{lev}
Let $\bm r=(r_1,\dots,r_n)$ and $\bm h=(h_1,\dots,h_m)$ be sequences over an alphabet. An \emph{edit} is one substitution
($r_i\to h_j$), deletion (drop $r_i$) or insertion (add $h_j$). The Levenshtein distance $d(\bm r,\bm h)$ is the minimum number of
edits that turns $\bm r$ into $\bm h$.
\end{definition}

\begin{theorem}{Dynamic programme for edit distance}{levdp}
Let $D[i,j]=d\big((r_1..r_i),(h_1..h_j)\big)$. Then $D[i,0]=i$, $D[0,j]=j$ and for $i,j\ge1$
\[
D[i,j]=\min\Big\{\,D[i-1,j]+1,\ \ D[i,j-1]+1,\ \ D[i-1,j-1]+\1[r_i\ne h_j]\,\Big\}.
\]
$d(\bm r,\bm h)=D[n,m]$, computed in $O(nm)$ time.
\end{theorem}
\begin{proof}
An optimal edit script, read left to right, aligns the prefixes. Look at its last operation. Either $r_i$ was deleted (cost 1 plus
the best way to turn $r_{1..i-1}$ into $h_{1..j}$), or $h_j$ was inserted (cost 1 plus $D[i,j-1]$), or $r_i$ was aligned with $h_j$
(cost 0 if equal, 1 if substituted, plus $D[i-1,j-1]$). Every script ends in one of these three ways, and each sub-problem is
solved optimally by induction on $i+j$, so the minimum of the three is optimal. The base cases need $i$ deletions or $j$ insertions.
\end{proof}

\begin{example}{``kitten'' $\to$ ``sitting''}{kitten}
\[
\begin{array}{c|cccccccc}
 & \epsilon & s & i & t & t & i & n & g\\\hline
\epsilon & 0&1&2&3&4&5&6&7\\
k & 1&\bm1&2&3&4&5&6&7\\
i & 2&2&\bm1&2&3&4&5&6\\
t & 3&3&2&\bm1&2&3&4&5\\
t & 4&4&3&2&\bm1&2&3&4\\
e & 5&5&4&3&2&\bm2&3&4\\
n & 6&6&5&4&3&3&\bm2&\bm3
\end{array}
\]
$d=3$: substitute $k\to s$, substitute $e\to i$, insert $g$. So $S=2$, $D=0$, $I=1$.
\end{example}

\section{WER and CER}

\begin{definition}{Word and character error rates}{wercer}
For a test set of utterances $i=1,\dots,K$ with references $\bm r^{(i)}$ and hypotheses $\bm h^{(i)}$, both \emph{normalised}
(Section~\ref{sec:norm}),
\[
\CER=\frac{\sum_{i} d_{\text{char}}\big(\bm r^{(i)},\bm h^{(i)}\big)}{\sum_i |\bm r^{(i)}|_{\text{char}}}
=\frac{S+D+I}{N},
\]
where sequences are compared character by character (spaces included). WER is the same with words (split on spaces) as symbols.
\end{definition}

Note the \emph{micro-average}: total errors over total reference length, not the mean of per-utterance rates. Long utterances weigh
more. This is what \texttt{jiwer.cer(list, list)} computes and what \file{metrics.py} reports.

\begin{proposition}{CER is not bounded by 1}{cerbound}
For one pair, $|m-n|\le d(\bm r,\bm h)\le\max(m,n)$. Hence $\CER\le\max(m,n)/n$, which is unbounded as $m\to\infty$.
\end{proposition}
\begin{proof}
Each edit changes the length by at most one, and the length must change by $|m-n|$, so $d\ge|m-n|$. Upper bound: substitute
$\min(m,n)$ aligned positions and insert or delete the remaining $|m-n|$ symbols, $\min(m,n)+|m-n|=\max(m,n)$ edits.
\end{proof}

\begin{intuition}[: what CER $>1$ means]
A CER above 1 means the hypothesis is much \emph{longer} than the reference. In this project it almost always means a
\emph{repetition loop}: the decoder keeps emitting the same phrase until it hits the 440-token limit. Vanilla Whisper's CER of
2.60 on IndicVoices Assamese is a loop symptom, not a ``very precise'' measure of quality.
\end{intuition}

\paragraph{Why CER is the primary metric here.}
\begin{itemize}
\item Indic words are long and inflected (Telugu is agglutinative): one wrong suffix makes the whole word wrong for WER, while CER
gives partial credit.
\item Spelling variation (e.g.\ with or without a nukta, two ways of writing a nasal) is common; CER penalises it by one character,
WER by a whole word.
\item Code-switched references write English words sometimes in Latin, sometimes in Indic script (Section~\ref{sec:csmetric}).
\end{itemize}
WER is still reported for completeness.

\section{Text normalisation}
\label{sec:norm}

Two strings that a human would call ``the same transcript'' can differ in Unicode bytes, punctuation and case. Scoring raw strings
would count these as errors. A \emph{single} normaliser is therefore applied to both the training targets and, at scoring time, to
both references and hypotheses.

\begin{definition}{The project's normaliser \file{normalize(text, lang)}}{normalizer}
\begin{enumerate}
\item Unicode NFC (canonical composition).
\item Delete bracketed spans: \texttt{<...>}, \texttt{[...]}, \texttt{(...)} (tags such as \texttt{<unintelligible>}).
\item Remove zero-width characters (ZWJ, ZWNJ, ZWSP, BOM).
\item For Indic languages: the IndicNLP normaliser of that language (canonical forms of nukta letters, chandrabindu, etc.).
\item For English: NFKC; keep in-word apostrophes, mapping the curly apostrophe U+2019 to \texttt{'}.
\item Lower-case.
\item Replace every Unicode \emph{symbol} (category S*) or \emph{punctuation} (category P*) character by a space. This includes the
danda \deva{।} (U+0964). Combining marks (category M*: matras, viramas) are \emph{kept}.
\item Collapse runs of whitespace; strip.
\end{enumerate}
Digits are \emph{not} converted to words.
\end{definition}

\begin{example}{Why NFC matters}{nfc}
The Devanagari letter \deva{क़} (qa) can be stored as one code point U+0958 or as two, \deva{क} U+0915 followed by the nukta U+093C.
They look identical, but a character-level comparison would count one substitution plus one insertion. NFC (plus IndicNLP's nukta
handling) maps both to the same sequence.
\end{example}

\begin{vivatip}[: two known side effects of the normaliser]
(1) Digits are not verbalised, so ``25'' vs ``twenty five'' counts as errors; this hurts Svarah (Chapter~\ref{ch:svarah}).
(2) Parenthesised text is deleted. For IndicVoices this removes tags, but in FLEURS the parenthesised words are usually \emph{read
aloud}, so 5--8\% of FLEURS rows lose spoken words from both targets and references. Re-scoring shows this inflates FLEURS CER by
0.005--0.02 for \emph{every} system (e.g.\ Stage 2 English 0.042 $\to$ 0.028 if kept). Rankings barely change. This was found in
the fourth code review and left open because fixing it changes the training targets.
\end{vivatip}

\section{LID metrics}

\begin{definition}{LID accuracy and confusion matrix}{lidacc}
For a test set whose true language is $\ell$, LID accuracy is the fraction of utterances with $\hat\ell=\ell$. Pooling all test sets,
the confusion matrix $C\in\mathbb N^{7\times7}$ has $C_{\ell\ell'}=\#\{\text{utterances of language }\ell\text{ predicted as }\ell'\}$.
Recall of $\ell$ is $C_{\ell\ell}/\sum_{\ell'}C_{\ell\ell'}$. \emph{Macro} accuracy averages per-set accuracies with equal weight.
\end{definition}

\section{Oracle versus predicted language}
\label{sec:oracle}

Whisper must be given a language token before it transcribes. We decode every test utterance twice:
\begin{itemize}
\item \textbf{Oracle LID}: the true language token is forced. This measures pure recognition quality.
\item \textbf{Predicted LID}: the model's own LID decision is used. This is what a user experiences.
\end{itemize}
The difference $\CER_{\text{pred}}-\CER_{\text{oracle}}$ is the \emph{routing cost} of LID errors. The literature reports about
4 CER points for a strong LID-then-ASR system (TalTech, ML-SUPERB 2.0); we measure ours in Chapter~\ref{ch:results}.

\begin{incode}
\file{eval_whisper.py} decodes with the oracle token, predicts the language, and \emph{re-decodes only the utterances whose
predicted language differs} (the same token gives the same greedy decode, so the oracle output is reused). It reports
\texttt{oracle\_norm}, \texttt{oracle\_raw}, \texttt{pred\_norm}, \texttt{lid\_acc}, \texttt{lid\_conf}, and the LID-head accuracy.
\end{incode}

\begin{exercise}{Micro vs macro}{micro}
Set A has 2 utterances: reference lengths 10 and 90 characters, with 5 and 9 errors. Compute the micro CER and the mean of per-utterance CERs.\par
\textit{Solution.} Micro: $(5+9)/(10+90)=0.14$. Mean of rates: $(0.5+0.1)/2=0.30$. The short utterance dominates the macro value.
The project uses micro CER within a test set, and reports the plain mean over test sets as the summary (``mean over 23 sets'').
\end{exercise}
